\documentclass[11pt,a4paper]{article}

\usepackage[utf8]{inputenc}
\usepackage[T1]{fontenc}
\usepackage{lmodern}
\usepackage{amsmath,amssymb,mathtools}
\usepackage[margin=2.5cm]{geometry}
\usepackage{hyperref}
\usepackage{booktabs}
\usepackage{xcolor}
\usepackage{fancyhdr}
\usepackage{enumitem}

\definecolor{statusamber}{RGB}{155,96,0}
\definecolor{statusblue}{RGB}{0,72,130}
\definecolor{statusgreen}{RGB}{0,120,50}

\pagestyle{fancy}
\fancyhf{}
\fancyhead[L]{\small\textsc{SCT Theory --- Prediction Card}}
\fancyhead[R]{\small\textsc{MT-1}}
\fancyfoot[C]{\thepage}

\hypersetup{
  colorlinks=true,
  linkcolor=blue!60!black,
  citecolor=green!40!black,
  urlcolor=blue!60!black,
  pdftitle={SCT Prediction Card: Schwarzschild Black Hole Entropy},
  pdfauthor={David Alfyorov},
}

\title{%
  \textbf{Prediction Card:}\\[4pt]
  \Large Schwarzschild Black Hole Entropy in SCT Theory
}
\author{David Alfyorov}
\date{April 2026 (v2)}

\begin{document}
\maketitle

\begin{center}
\end{center}

\begin{center}
  \colorbox{statusgreen!15}{\parbox{0.82\textwidth}{\centering
    \textbf{\color{statusgreen} STATUS: COMPLETE (CONDITIONAL on second law)}\\[2pt]
    \small Schwarzschild BH entropy: \textbf{CERTIFIED}.\\
    174+ checks (154 pytest/verification $+$ 14 Lean theorems $+$ 6 D-R checks).\\[2pt]
    6-agent pipeline (L$\to$L-R$\to$D$\to$D-R$\to$V$\to$V-R): ALL PASS.\\[4pt]
    \textit{Second law:} \textbf{CONDITIONAL} on MR-2 ghost resolution (Wall stability).\\
    \textit{Kerr generalization:} DEFERRED (same topological correction expected).\\
    \textit{Greybody:} Schematic approximation only (not rigorous Regge--Wheeler).
  }}
\end{center}

%% ==========================================================
\section{Source}
%% ==========================================================

\begin{itemize}[leftmargin=2em]
  \item \textbf{Immediate inputs:}
  \begin{itemize}
    \item NT-1b Phase~3: SM-summed form factor coefficients
    $\alpha_C = 13/120$ (parameter-free), $\alpha_R(\xi) = 2(\xi-1/6)^2$.
    \item NT-2: Entire-function property of $F_1(z)$, $F_2(z,\xi)$ (no spurious poles).
    \item NT-4b: Full nonlinear field equations and the Wald entropy variation
    $\delta S_{\mathrm{spec}}/\delta R_{\mu\nu\rho\sigma}$.
    \item MR-2: Ghost catalogue and unitarity conditions (conditional input for second law).
  \end{itemize}
  \item \textbf{Key references:}
    Wald (1993), gr-qc/9307038;
    Iyer--Wald (1994), gr-qc/9403028;
    Jacobson--Myers (1993), PRL~\textbf{70}, 3684;
    Sen (2012), arXiv:\texttt{1205.0971};
    Wall (2015), arXiv:\texttt{1504.08040};
    Bekenstein (1973), PRD~\textbf{7}, 2333;
    Hawking (1975), CMP~\textbf{43}, 199.
\end{itemize}

%% ==========================================================
\section{Observable Quantity}
%% ==========================================================

The observable is the entropy of a Schwarzschild black hole of mass~$M$,
including quantum corrections to the classical Bekenstein--Hawking area law:
\[
  S_{\mathrm{BH}} = S_{\mathrm{classical}} + S_{\mathrm{quantum}}.
\]
The quantum corrections encode the ultraviolet structure of the gravitational
theory and can discriminate between competing approaches to quantum gravity.

%% ==========================================================
\section{SCT Prediction}
%% ==========================================================

\begin{equation}
\boxed{
  S = \frac{A}{4G}
  + \frac{13}{120\pi}
  + \frac{37}{24}\,\ln\!\Bigl(\frac{A}{\ell_P^2}\Bigr)
  + \mathcal{O}(1)
}
\label{eq:main}
\end{equation}

where $A = 16\pi G^2 M^2/c^4$ is the horizon area and $\ell_P^2 = \hbar G/c^3$.

\subsection{Term-by-term decomposition}

\begin{enumerate}
  \item \textbf{Bekenstein--Hawking term:} $A/(4G)$.\\
  Standard area law from the Einstein--Hilbert sector.
  Recovery of classical GR in the $\Lambda \to \infty$ limit.

  \item \textbf{Topological Weyl correction:}
  $\Delta S_{\mathrm{Weyl}} = \alpha_C/\pi = 13/(120\pi) \approx 0.03448$.\\
  Computed via the Jacobson--Myers formula applied to the $C^2$ sector of the spectral action.
  Depends on the Euler characteristic $\chi(\mathcal{H}) = 2$ (for $S^2$ horizon topology)
  but \emph{not} on the horizon area --- a topological invariant.

  \item \textbf{One-loop logarithmic correction:}
  $c_{\log} = 37/24 \approx 1.542$.\\
  Computed from the Sen (2012) formula:
  \[
    c_{\log} = \frac{1}{180}\bigl[2 N_s + 7 N_D - 26 N_V + 424\bigr],
  \]
  with SM field content $N_s = 4$, $N_D = 22.5$, $N_V = 12$:
  \[
    c_{\log} = \frac{8 + 157.5 - 312 + 424}{180}
    = \frac{277.5}{180} = \frac{37}{24}.
  \]
  \textbf{Parameter-free}: determined entirely by the Standard Model field content
  plus the graviton sector. No dependence on the spectral cutoff~$\Lambda$
  or the non-minimal coupling~$\xi$.

  \item \textbf{$R^2$ sector:} Vanishes identically on Schwarzschild ($R = 0$).
\end{enumerate}

\subsection{Hierarchy}

For astrophysical black holes:
\[
  S_{\mathrm{BH}} \gg \Delta S_{\log} \gg \Delta S_{\mathrm{Weyl}}.
\]

\begin{center}
\begin{tabular}{lccc}
\toprule
$M/M_\odot$ & $S_{\mathrm{BH}}$ & $\Delta S_{\log}$ & $\Delta S_{\mathrm{Weyl}}$ \\
\midrule
1         & $1.05 \times 10^{77}$ & 275.5 & 0.0345 \\
10        & $1.05 \times 10^{79}$ & 282.6 & 0.0345 \\
100       & $1.05 \times 10^{81}$ & 289.7 & 0.0345 \\
$10^6$    & $1.05 \times 10^{89}$ & 318.1 & 0.0345 \\
$4.3\times10^6$ (Sgr~A*) & $1.94 \times 10^{90}$ & 322.6 & 0.0345 \\
\bottomrule
\end{tabular}
\end{center}

%% ==========================================================
\section{Key Discriminant: $c_{\log}$ Sign}
%% ==========================================================

\begin{center}
\begin{tabular}{lrl}
\toprule
\textbf{Approach} & $c_{\log}$ & \textbf{Notes} \\
\midrule
\textbf{SCT (this work)} & $+37/24 \approx +1.542$ & Parameter-free; SM + graviton via Sen \\
LQG (isolated horizons) & $-3/2 = -1.500$ & Opposite sign (Kaul--Majumdar 2000) \\
String theory & model-dependent & Matches Sen for specific extremal BH \\
Asymptotic safety & not well-defined & Running $G(k)$ complicates definition \\
IDG (Modesto et al.) & same Sen formula & Same field content $\Rightarrow$ identical $c_{\log}$ \\
\bottomrule
\end{tabular}
\end{center}

The SCT prediction $c_{\log} = +37/24$ has the \textbf{opposite sign} from the LQG
result $c_{\log}^{\mathrm{LQG}} = -3/2$. In~SCT, quantum corrections \emph{increase}
entropy beyond the area law; in~LQG they \emph{decrease} it. This qualitative
distinction is a potential observational discriminant if logarithmic corrections
to black hole entropy become accessible.

%% ==========================================================
\section{Four Laws of Black Hole Thermodynamics}
%% ==========================================================

\begin{center}
\begin{tabular}{lll}
\toprule
\textbf{Law} & \textbf{Status} & \textbf{Key Evidence} \\
\midrule
Zeroth & \textcolor{statusgreen}{PASS} & $\kappa = \mathrm{const}$ by $S^2$ symmetry \\
First  & \textcolor{statusgreen}{PASS} & Iyer--Wald theorem; numerical $\delta = 6.2\times10^{-6}$ \\
Second & \textcolor{statusamber}{CONDITIONAL} & Depends on MR-2 ghost resolution (Wall stability) \\
Third  & \textcolor{statusgreen}{PASS} & $\kappa \to 0$ only as $M \to \infty$; SCT corrections $\sim 10^{-62}$ \\
\bottomrule
\end{tabular}
\end{center}

%% ==========================================================
\section{Modified Hawking Radiation}
%% ==========================================================

The Hawking temperature $T_H = \hbar c^3/(8\pi G M k_B)$ is unmodified at leading
order. The nonlocal form factors modify greybody factors at the scale $T_H/\Lambda$:
\[
  \Gamma^{\mathrm{SCT}}(\omega) = \Gamma^{\mathrm{GR}}(\omega)\cdot\bigl[1 + \mathcal{O}(\omega^2/\Lambda^2)\bigr].
\]
For all astrophysical black holes, $T_H/\Lambda \sim 10^{-10}$ (at $\Lambda = 2.38$~meV),
making SCT modifications negligible.

\textbf{Caution:} The greybody estimate is schematic (order-of-magnitude),
not a rigorous Regge--Wheeler calculation. Precision QNM predictions require
the downstream LT-3a analysis.

%% ==========================================================
\section{Verification}
%% ==========================================================

\begin{center}
\begin{tabular}{lr}
\toprule
\textbf{Verification source} & \textbf{Count} \\
\midrule
pytest regression tests & 82 \\
8-layer independent verification & 146 \\
Lean~4 theorems (sorry-free) & 14 \\
D-R independent re-derivation checks & 6 \\
V-R independent spot checks & 8 \\
\midrule
\textbf{Total} & $\mathbf{174+}$ \\
\bottomrule
\end{tabular}
\end{center}

Dual derivation: Jacobson--Myers formula vs.\ direct Wald Noether charge
integration on the Schwarzschild horizon. Both yield $\Delta S = 13/(120\pi)$.

%% ==========================================================
\section{Dependencies}
%% ==========================================================

\begin{itemize}[leftmargin=2em]
  \item \textbf{Required (satisfied):} NT-1b Phase~3 ($\alpha_C = 13/120$),
  NT-2 (entire-function property), NT-4b (nonlinear field equations).
  \item \textbf{Conditional:} MR-2 (ghost resolution, for second law only).
  \item \textbf{Downstream:} MR-9 (singularity resolution),
  LT-2 (information paradox), LT-3a (QNM predictions),
  COMP-1 (cross-program comparison).
\end{itemize}

%% ==========================================================
\section{Limitations}
%% ==========================================================

\begin{enumerate}
  \item \textbf{Second law is CONDITIONAL} on MR-2 ghost resolution
  (Wall's linearized stability condition). The ghost poles in
  $\Pi_{\mathrm{TT}}(z)$ require fakeon prescription resolution.
  \item \textbf{Greybody factors are schematic} ---
  not a rigorous Regge--Wheeler calculation.
  \item \textbf{Schwarzschild only} --- Kerr generalization not yet derived
  (same topological correction expected, different logarithmic coefficient).
  \item \textbf{Nonlocal Wald entropy} --- the local-limit approximation
  $F_1(0) = 1$ may not hold for Planck-scale black holes.
  \item \textbf{Graviton zero modes} --- subtleties in the zero-mode
  sector of Sen's formula could modify the $+424$ contribution
  at extreme precision.
\end{enumerate}

\end{document}
